\chapter{2012年江西卷（理）}

\section{单选题}

\begin{problem}
若集合 \(A = \{-1, 1\}\)，\(B = \{0, 2\}\)，则集合 \(\{z \mid z = x + y, x \in A, y \in B\}\) 中的元素个数为（\quad）

\choices
    {\(5\)}
    {\(4\)}
    {\(3\)}
    {\(2\)}
\end{problem}

\begin{answer}
C
\end{answer}

\begin{solution}
由 \(x\in A\)、\(y\in B\)，可得
\(x+y\in\{-1+0,-1+2,1+0,1+2\}=\{-1,1,3\}\).
集合中的不同元素只有 \(3\) 个，所以应选 C.
\end{solution}

\begin{problem}
下列函数中，与函数 \( y = \frac{1}{\sqrt[3]{x}} \) 定义域相同的函数为（\quad）

\choices
    {\( y = \frac{1}{\sin x} \)}
    {\( y = \frac{\ln x}{x} \)}
    {\( y = x\e^{x} \)}
    {\( y = \frac{\sin x}{x} \)}
\end{problem}

\begin{answer}
D
\end{answer}

\begin{solution}
函数 \(y=\dfrac{1}{\sqrt[3]{x}}\) 的分母不能为零，故其定义域为
\(\{x\mid x\neq0\}\).
选项 A 的定义域要求 \(\sin x\neq0\)，选项 B 要求 \(x>0\)，选项 C 的定义域为 \(\R\)，而选项 D 的定义域为 \(\{x\mid x\neq0\}\). 因此应选 D.
\end{solution}

\begin{problem}
若函数 \( f(x) = \begin{cases}
x^2 + 1, & x \leqslant 1,\\
\lg x, & x > 1,
\end{cases} \) 则 \( f(f(10)) = \)（\quad）

\choices
    {\( \lg 101 \)}
    {\(2\)}
    {\(1\)}
    {\(0\)}
\end{problem}

\begin{answer}
B
\end{answer}

\begin{solution}
因为 \(10>1\)，所以 \(f(10)=\lg10=1\). 又因为 \(1\leqslant1\)，代入第一段解析式得 \(f(1)=1^2+1=2\). 因而 \(f(f(10))=2\)，应选 B.
\end{solution}

\begin{problem}
若 \(\tan \theta + \frac{1}{\tan \theta} = 4\)，则 \(\sin 2\theta =\)（\quad）

\choices
    {\( \frac{1}{8} \)}
    {\( \frac{1}{4} \)}
    {\( \frac{1}{3} \)}
    {\( \frac{1}{2} \)}
\end{problem}

\begin{answer}
D
\end{answer}

\begin{solution}
令 \(t=\tan\theta\). 由题意 \(t\neq0\)，且
\(t+\dfrac1t=4\)，从而 \(t^2+1=4t\). 因此
\(\sin2\theta=\dfrac{2\tan\theta}{1+\tan^2\theta}=\dfrac{2t}{4t}=\dfrac12\)，应选 D.
\end{solution}

\begin{problem}
下列命题中，假命题为（\quad）

\choices
    {存在四边相等的四边形不是正方形}
    {\(z_{1}\)、\(z_{2}\in\mathbf{C}\)，\(z_{1}+z_{2}\) 为实数的充分必要条件是 \(z_{1}\)、\(z_{2}\) 互为共轭复数}
    {若 \(x\)、\(y\in\R\)，且 \(x+y>2\)，则 \(x\)、\(y\) 至少有一个大于 \(1\)}
    {对于任意 \(n\in\mathbf{N}_{+}\)，\(\mathrm C_{n}^{0}+\mathrm C_{n}^{1}+\cdots+\mathrm C_{n}^{n}\) 都是偶数}
\end{problem}

\begin{answer}
B
\end{answer}

\begin{solution}
选项 A 正确，例如菱形只要不是正方形，就存在四边相等而不是正方形的四边形.

选项 B 不正确. 取 \(z_1=1\)、\(z_2=2\)，则 \(z_1+z_2=3\) 是实数，但 \(z_1\)、\(z_2\) 不互为共轭复数，所以所述充分必要条件不成立.

选项 C 正确. 若 \(x\leqslant1\) 且 \(y\leqslant1\)，则 \(x+y\leqslant2\)，与 \(x+y>2\) 矛盾，故 \(x\)、\(y\) 至少有一个大于 \(1\).

选项 D 正确，因为二项式系数满足
\(\mathrm C_n^0+\mathrm C_n^1+\cdots+\mathrm C_n^n=2^n\)，而 \(n\in\N_+\) 时 \(2^n\) 为偶数. 所以假命题为 B.
\end{solution}

\begin{problem}
观察下列各式： \(a + b = 1\)，\(a^2 + b^2 = 3\)，\(a^3 + b^3 = 4\)，\(a^4 + b^4 = 7\)，\(a^5 + b^5 = 11\)，\(\cdots\)，则 \(a^{10} + b^{10} =\) （\quad）

\choices
    {\(28\)}
    {\(76\)}
    {\(123\)}
    {\(199\)}
\end{problem}

\begin{answer}
C
\end{answer}

\begin{solution}
记 \(s_n=a^n+b^n\)，并令 \(s_0=2\). 由 \(a+b=1\) 和 \(a^2+b^2=3\)，有
\(1=(a+b)^2=a^2+2ab+b^2=3+2ab\)，故 \(ab=-1\).

于是对 \(n\geqslant2\)，
\(s_n=(a+b)s_{n-1}-ab s_{n-2}=s_{n-1}+s_{n-2}\).
由 \(s_1=1\)、\(s_2=3\) 递推得
\(s_3=4\)、\(s_4=7\)、\(s_5=11\)、\(s_6=18\)、\(s_7=29\)、\(s_8=47\)、\(s_9=76\)、\(s_{10}=123\).
所以 \(a^{10}+b^{10}=123\)，应选 C.
\end{solution}

\begin{problem}
在直角三角形 \(ABC\) 中，点 \(D\) 是斜边 \(AB\) 的中点，点 \(P\) 为线段 \(CD\) 的中点，则 \(\frac{|PA|^2 + |PB|^2}{|PC|^2} =\)（\quad）

\choices
    {\(2\)}
    {\(4\)}
    {\(5\)}
    {\(10\)}
\end{problem}

\begin{answer}
D
\end{answer}

\begin{solution}
不妨在直角坐标系中取 \(C=(0,0)\)、\(A=(u,0)\)、\(B=(0,v)\)，其中 \(u\)，\(v>0\). 则
\(D=\left(\dfrac u2,\dfrac v2\right)\)、\(P=\left(\dfrac u4,\dfrac v4\right)\).

因此
\(\lvert PA\rvert^2=\dfrac{9u^2+v^2}{16}\)、\(\lvert PB\rvert^2=\dfrac{u^2+9v^2}{16}\)，而
\(\lvert PC\rvert^2=\dfrac{u^2+v^2}{16}\). 从而
\[
\frac{\lvert PA\rvert^2+\lvert PB\rvert^2}{\lvert PC\rvert^2}
=\frac{10(u^2+v^2)/16}{(u^2+v^2)/16}=10
\]
故应选 D.
\end{solution}

\begin{problem}
某农户计划种植黄瓜和韭菜，种植面积不超过 \(50\text{ 亩}\)，投入资金不超过 \(54\text{ 万元}\)，假设种植黄瓜和韭菜的产量、成本和售价如下表

\begin{center}
\begin{tabular}{c|c|c|c}
\hline
 & 年产量/亩 & 年种植成本/亩 & 每吨售价 \\
\hline
黄瓜 & \(4\text{ 吨}\) & \(1.2\text{ 万元}\) & \(0.55\text{ 万元}\) \\
\hline
韭菜 & \(6\text{ 吨}\) & \(0.9\text{ 万元}\) & \(0.3\text{ 万元}\) \\
\hline
\end{tabular}
\end{center}

为使一年的种植总利润（\(\text{总利润} = \text{总销售收入} - \text{总种植成本}\)）最大，那么黄瓜和韭菜的种植面积（单位：亩）分别为（\quad）

\choices
    {\(50,0\)}
    {\(30,20\)}
    {\(20,30\)}
    {\(0,50\)}
\end{problem}

\begin{answer}
B
\end{answer}

\begin{solution}
设黄瓜和韭菜的种植面积分别为 \(x\)、\(y\) 亩，则
\(x\geqslant0\)、\(y\geqslant0\)、\(x+y\leqslant50\)，且 \(1.2x+0.9y\leqslant54\).

每亩黄瓜的利润为 \(4\times0.55-1.2=1\) 万元，每亩韭菜的利润为 \(6\times0.3-0.9=0.9\) 万元，故总利润为 \(z=x+0.9y\). 可行域的有关顶点为
\((0,0)\)、\((45,0)\)、\((30,20)\)、\((0,50)\).

在这些顶点处，\(z\) 依次为 \(0\)、\(45\)、\(48\)、\(45\). 最大值在 \((30,20)\) 处取得，所以黄瓜和韭菜的种植面积分别为 \(30\) 亩和 \(20\) 亩，应选 B.
\end{solution}

\begin{problem}
样本 \((x_{1}, x_{2}, \cdots, x_{n})\) 的平均数为 \(\overline{x}\)，样本 \((y_{1}, y_{2}, \cdots, y_{m})\) 的平均数为 \(\overline{y}\) （记 \(\overline{x} \neq \overline{y}\)）. 若样本 \((x_{1}, x_{2}, \cdots, x_{n}, y_{1}, y_{2}, \cdots, y_{m})\) 的平均数 \(\overline{z} = \alpha \overline{x} + (1 - \alpha)\overline{y}\)，其中 \(0 < \alpha < \frac{1}{2}\)，则 \(n\)，\(m\) 的大小关系为（\quad）

\choices
    {\(n < m\)}
    {\(n > m\)}
    {\(n = m\)}
    {不能确定}
\end{problem}

\begin{answer}
A
\end{answer}

\begin{solution}
合并样本的平均数满足
\(\overline z=\dfrac{n\overline x+m\overline y}{n+m}\).
又因 \(\overline x\neq\overline y\)，题设等式中两种平均数的系数必须分别相等，因此
\(\dfrac n{n+m}=\alpha\).
由 \(0<\alpha<\dfrac12\) 得 \(n<m\)，所以应选 A.
\end{solution}

\begin{problem}
如图，已知正四棱锥 \(S - ABCD\) 所有棱长都为 1，点 \(E\) 是侧棱 \(SC\) 上一动点，过点 \(E\) 垂直于 \(SC\) 的截面将正四棱锥分成上、下两部分. 记 \(SE = x\) （\(0 < x < 1\)），截面下面部分的体积为 \(V(x)\)，则函数 \(y = V(x)\) 的图象大致为（\quad）

\bitmapfigure[width=0.42\linewidth]{img/2012/jiangxi_science/q10_fig2.png}

\choices
    {\choicebitmap[width=0.15\paperwidth]{img/2012/jiangxi_science/q10_optA.png}}
    {\choicebitmap[width=0.15\paperwidth]{img/2012/jiangxi_science/q10_optB.png}}
    {\choicebitmap[width=0.15\paperwidth]{img/2012/jiangxi_science/q10_optC.png}}
    {\choicebitmap[width=0.15\paperwidth]{img/2012/jiangxi_science/q10_optD.png}}
\end{problem}

\begin{answer}
A
\end{answer}

\begin{solution}
设正四棱锥的底面为边长 \(1\) 的正方形，高为 \(h\). 由顶点到底面中心的距离为 \(\dfrac{1}{\sqrt2}\)，且侧棱长为 \(1\)，得 \(h=\dfrac1{\sqrt2}\)，故整个棱锥体积为 \(\dfrac{\sqrt2}{6}\).

取底面中心为原点，令 \(C=\left(\dfrac12,\dfrac12,0\right)\)、\(S=\left(0,0,\dfrac1{\sqrt2}\right)\)，并令 \(E=S+x(C-S)\). 截面平面方程为
\[
\frac12(X+Y+1-2x)-\frac{Z}{\sqrt2}=0
\]
将截面上的 \(Z\) 写成 \(Z=(X+Y+1-2x)/\sqrt2\). 正四棱锥的侧面在底面投影 \((X,Y)\) 上满足
\(-\frac12\leqslant X\)，\(Y\leqslant\frac12\)，\(0\leqslant X+Y+1-2x\leqslant1-2|X|\)，\(0\leqslant X+Y+1-2x\leqslant1-2|Y|\).
记 \(u=X+Y\)、\(v=X-Y\). 因为 \(X=(u+v)/2\)、\(Y=(u-v)/2\)，所以
\(\max\{|X|,|Y|\}=\frac{|u|+|v|}{2}\)，\(\mathrm dX\,\mathrm dY=\frac12\mathrm du\,\mathrm dv\).

正四棱锥在底面投影上的上表面高度为 \((1-2\max\{|X|,|Y|\})/\sqrt2\). 因而截面投影区域满足
\(u+1-2x\geqslant0\)，\(u+|u|+|v|\leqslant2x\)，\(|u|+|v|\leqslant1\).

当 \(0\leqslant x\leqslant\dfrac12\) 时，若 \(2x-1\leqslant u\leqslant0\)，则 \(1+u\geqslant2x\)，所以 \(|v|\leqslant2x\)；若 \(0\leqslant u\leqslant x\)，则 \(2(x-u)\leqslant1-u\)，所以 \(|v|\leqslant2(x-u)\). 因此投影面积为
\[
A_0(x)=\frac12\left(\int_{2x-1}^{0}4x\,\mathrm du+\int_{0}^{x}4(x-u)\,\mathrm du\right)=x(2-3x)
\]

当 \(\dfrac12\leqslant x\leqslant1\) 时，\(2x-1\leqslant u\leqslant x\)，且 \(2(x-u)\leqslant1-u\)，故 \(|v|\leqslant2(x-u)\)，从而
\[
A_0(x)=\frac12\int_{2x-1}^{x}4(x-u)\,\mathrm du=(1-x)^2
\]

截面面积为 \(\sqrt2 A_0(x)\)，且 \(V'(x)\) 等于截面面积的相反数. 由 \(V(0)=\dfrac{\sqrt2}{6}\) 积分得
\[
V(x)=\begin{cases}
\sqrt2\left(\dfrac16-x^2+x^3\right),&0\leqslant x\leqslant\dfrac12,\\
\dfrac{\sqrt2}{3}(1-x)^3,&\dfrac12\leqslant x\leqslant1.
\end{cases}
\]
该函数从 \(\dfrac{\sqrt2}{6}\) 单调下降到 \(0\). 在 \(0<x<\dfrac13\) 时 \(V''(x)<0\)，在 \(\dfrac13<x<1\) 时 \(V''(x)>0\)，且 \(x=\dfrac12\) 处函数及其一、二阶导数连续，图象与选项 A 相符，所以应选 A.
\end{solution}

\section{填空题}

\begin{problem}
计算定积分 \(\int_{-1}^{1}\left(x^{2} + \sin x\right)\mathrm{d}x =\fillinblank{}\).
\end{problem}

\begin{answer}
\(\frac{2}{3}\)
\end{answer}

\begin{solution}
在区间 \([-1,1]\) 上，\(x^2\) 是偶函数，\(\sin x\) 是奇函数，因此
\[
\int_{-1}^{1}\left(x^2+\sin x\right)\mathrm{d}x
=2\int_0^1x^2\mathrm{d}x+\int_{-1}^{1}\sin x\,\mathrm{d}x
=2\left[\frac{x^3}{3}\right]_0^1+0
=\frac23
\]
\end{solution}

\begin{problem}
设数列 \(\{a_{n}\}\)，\(\{b_{n}\}\) 都是等差数列，若 \(a_{1} + b_{1} = 7\)，\(a_{3} + b_{3} = 21\)，则 \(a_{5} + b_{5} =\fillinblank{}\).
\end{problem}

\begin{answer}
35
\end{answer}

\begin{solution}
令 \(c_n=a_n+b_n\). 两个等差数列对应项相加仍为等差数列，设其公差为 \(d\). 由
\(c_1=7\)，\(c_3=21\)
得 \(2d=c_3-c_1=14\)，所以 \(d=7\). 因而
\[
c_5=c_1+4d=7+4\times7=35
\]
所以 \(a_5+b_5=35\).
\end{solution}

\begin{problem}
椭圆 \(\frac{x^{2}}{a^{2}} + \frac{y^{2}}{b^{2}} = 1\)（\(a > b > 0\)）的左、右顶点分别是 \(A\)，\(B\)，左、右焦点分别是 \(F_1\)，\(F_2\). 若 \(|AF_1|\)，\(|F_1F_2|\)，\(|F_1B|\) 成等比数列，则此椭圆的离心率为\fillinblank{}.
\end{problem}

\begin{answer}
\(\frac{\sqrt{5}}{5}\)
\end{answer}

\begin{solution}
设椭圆的半焦距为 \(c\)，则
\(A=(-a,0)\)，\(B=(a,0)\)，\(F_1=(-c,0)\)，\(F_2=(c,0)\)
且 \(b^2=a^2-c^2\). 因此
\(|AF_1|=a-c\)，\(|F_1F_2|=2c\)，\(|F_1B|=a+c\).
三项成等比数列，故中间项的平方等于两端项的乘积：
\[
(2c)^2=(a-c)(a+c)=a^2-c^2=b^2
\]
所以 \(a^2=5c^2\)，离心率
\[
e=\frac ca=\frac1{\sqrt5}=\frac{\sqrt5}{5}
\]
\end{solution}

\begin{problem}
如图为某算法的程序框图，则程序运行后输出的结果是\fillinblank{}.

\bitmapfigure[max width=0.62\linewidth]{img/2012/jiangxi_science/q14_flowchart.png}
\end{problem}

\begin{answer}
3
\end{answer}

\begin{solution}
按程序框图逐次计算. 初始令 \(T=0\)，\(k=1\)，每次将
\(\sin\dfrac{k\pi}{2}-\sin\dfrac{(k-1)\pi}{2}\) 加到 \(T\) 中，并在 \(k=5\) 后结束循环.

当 \(k=1\)，\(2\)，\(3\)，\(4\)，\(5\) 时，所比较的两项分别为
\begin{center}
\begin{tabular}{c|ccccc}
\(k\) & \(1\) & \(2\) & \(3\) & \(4\) & \(5\) \\
\hline
\(\sin\dfrac{k\pi}{2}\) & \(1\) & \(0\) & \(-1\) & \(0\) & \(1\) \\
\(\sin\dfrac{(k-1)\pi}{2}\) & \(0\) & \(1\) & \(0\) & \(-1\) & \(0\)
\end{tabular}
\end{center}
因此不等式
\(\sin\dfrac{k\pi}{2}>\sin\dfrac{(k-1)\pi}{2}\)
在 \(k=1\)，\(4\)，\(5\) 时成立，在 \(k=2\)，\(3\) 时不成立. 满足条件的次数为 \(3\)，故程序输出 \(3\).
\end{solution}

\section{选做题：填空题}

\begin{problem}
曲线 \(C\) 的直角坐标方程为 \(x^{2} + y^{2} - 2x = 0\)，以原点为极点，\(x\) 轴的正半轴为极轴建立极坐标系，则曲线 \(C\) 的极坐标方程为\fillinblank{}.
\end{problem}

\begin{answer}
\(\rho = 2\cos\theta\)
\end{answer}

\begin{solution}
由极坐标与直角坐标的关系
\(x=\rho\cos\theta\)，\(x^2+y^2=\rho^2\)
将其代入 \(x^2+y^2-2x=0\)，得
\[
\rho^2-2\rho\cos\theta=0
\]
即 \(\rho(\rho-2\cos\theta)=0\). 原点本身在曲线 \(C\) 上，而极坐标方程 \(\rho=2\cos\theta\) 在 \(\cos\theta=0\) 时也包含原点，所以曲线的极坐标方程为
\[
\rho=2\cos\theta
\]
\end{solution}

\begin{problem}
在实数范围内，不等式 \(|2x - 1| + |2x + 1| \leqslant 6\) 的解集为\fillinblank{}.
\end{problem}

\begin{answer}
\(\left[-\frac{3}{2}, \frac{3}{2}\right]\)
\end{answer}

\begin{solution}
绝对值内的两个一次式分别在 \(x=\pm\dfrac12\) 处变号，分三段讨论.

当 \(x\leqslant-\dfrac12\) 时，
\[
|2x-1|+|2x+1|=(1-2x)+(-2x-1)=-4x
\]
不等式化为 \(-4x\leqslant6\)，即 \(x\geqslant-\dfrac32\)，与本段合并得
\(\left[-\dfrac32,-\dfrac12\right]\).

当 \(-\dfrac12\leqslant x\leqslant\dfrac12\) 时，
\[
|2x-1|+|2x+1|=(1-2x)+(2x+1)=2\leqslant6
\]
本段所有 \(x\) 都满足.

当 \(x\geqslant\dfrac12\) 时，
\[
|2x-1|+|2x+1|=(2x-1)+(2x+1)=4x
\]
不等式化为 \(4x\leqslant6\)，即 \(x\leqslant\dfrac32\)，与本段合并得
\(\left[\dfrac12,\dfrac32\right]\).

合并三段解集，得到
\[
\left[-\frac32,\frac32\right]
\]
\end{solution}

\section{解答题}

\begin{problem}
已知数列 \( \{a_{n}\} \) 的前 \(n\) 项和 \( S_{n} = -\frac{1}{2}n^{2} + kn \) （其中 \( k \in \N_{+} \) ），且 \( S_{n} \) 的最大值为 8.
\begin{enumerate}
\item 确定常数 \( k \)，并求 \( a_{n} \)；
\item 求数列 \( \left\{\frac{9-2a_{n}}{2^{n}}\right\} \) 的前 \(n\) 项和 \( T_{n} \).
\end{enumerate}
\end{problem}

\begin{solution}
\begin{enumerate}
\item 将 \(S_n\) 配方，得
\[
S_n=-\frac12(n-k)^2+\frac{k^2}{2}
\]
由于 \(n\in\mathbf N_+\)，且 \(k\in\mathbf N_+\)，当 \(n=k\) 时取得最大值. 因此
\(\frac{k^2}{2}=8\)，\(k=4\).
由前 \(n\) 项和与通项的关系，
\[
a_1=S_1=-\frac12+4=\frac72
\]
当 \(n\geqslant2\) 时，
\[
a_n=S_n-S_{n-1}
=\left(-\frac12n^2+4n\right)
-\left[-\frac12(n-1)^2+4(n-1)\right]
=\frac92-n
\]
\item 记
\[
b_n=\frac{9-2a_n}{2^n}
\]
由 \(a_1=\frac72\) 得 \(b_1=1\)；当 \(n\geqslant2\) 时，由 \(a_n=\frac92-n\) 得
\[
b_n=\frac{9-2\left(\frac92-n\right)}{2^n}
=\frac{n}{2^{n-1}}
\]
故 \(b_n=\dfrac{n}{2^{n-1}}\) 对 \(n=1\) 也成立，从而
\[
T_n=\sum_{i=1}^n\frac{i}{2^{i-1}}
\]
由有限等比数列求和公式
\[
\sum_{i=1}^n i r^{i-1}
=\left(\sum_{i=1}^n r^i\right)'
=\left(\frac{r-r^{n+1}}{1-r}\right)'
=\frac{1-(n+1)r^n+nr^{n+1}}{(1-r)^2}
\]
令 \(r=\frac12\)，可得
\[
T_n
=4\left(1-\frac{n+1}{2^n}+\frac{n}{2^{n+1}}\right)
=4-\frac{n+2}{2^{n-1}}
\]
\end{enumerate}
\end{solution}

\begin{problem}
在 \(\triangle ABC\) 中，角 \(A\)，\(B\)，\(C\) 的对边分别为 \(a\)，\(b\)，\(c\). 已知 \(A = \frac{\pi}{4}\)，\(b \sin \left(\frac{\pi}{4} + C\right) - c \sin \left(\frac{\pi}{4} + B\right) = a\).
\begin{enumerate}
\item 求证： \(B - C = \frac{\pi}{2}\)；
\item 若 \(a = \sqrt{2}\)，求 \(\triangle ABC\) 的面积.
\end{enumerate}
\end{problem}

\begin{solution}
\begin{enumerate}
\item 设外接圆半径为 \(R\). 由正弦定理，
\(a=2R\sin A\)，\(b=2R\sin B\)，\(c=2R\sin C\).
因 \(A=\dfrac\pi4\)，题设左端为
\[
b\sin\left(\frac\pi4+C\right)
-c\sin\left(\frac\pi4+B\right)
=\frac1{\sqrt2}\bigl[b(\cos C+\sin C)-c(\cos B+\sin B)\bigr]
\]
代入正弦定理，并利用
\(\sin B\cos C-\sin C\cos B=\sin(B-C)\)，上式化为
\[
\sqrt2R\bigl(\sin B\cos C-\sin C\cos B\bigr)
=\sqrt2R\sin(B-C)
\]
而 \(a=2R\sin\dfrac\pi4=\sqrt2R\)，所以题设等式给出
\[
\sin(B-C)=1
\]
三角形内角满足 \(B>0\)、\(C>0\)、\(B+C=\dfrac{3\pi}{4}\)，故
\[
-\frac{3\pi}{4}<B-C<\frac{3\pi}{4}
\]
在此范围内 \(\sin(B-C)=1\) 只能有
\[
B-C=\frac\pi2
\]
\item 由 \(B+C=\dfrac{3\pi}{4}\) 与 \(B-C=\dfrac\pi2\)，解得
\(B=\frac{5\pi}{8}\)，\(C=\frac\pi8\).
又 \(a=\sqrt2\)，由正弦定理
\[
2R=\frac{a}{\sin A}=\frac{\sqrt2}{\sqrt2/2}=2
\]
所以 \(R=1\). 于是
\(b=2\sin\frac{5\pi}{8}\)，\(c=2\sin\frac\pi8\).
三角形面积为
\[
\frac12bc\sin A
=\frac12\cdot4\sin\frac{5\pi}{8}\sin\frac\pi8\cdot\frac{\sqrt2}{2}
\]
由于 \(\sin\dfrac{5\pi}{8}=\sin\dfrac{3\pi}{8}\)，且
\[
\sin\frac{3\pi}{8}\sin\frac\pi8
=\frac{\cos\frac\pi4-\cos\frac\pi2}{2}
=\frac{\sqrt2}{4}
\]
故面积为
\[
\frac12\cdot4\cdot\frac{\sqrt2}{4}\cdot\frac{\sqrt2}{2}
=\frac12
\]
\end{enumerate}
\end{solution}

\begin{problem}
如图，从 \(A_{1}(1,0,0)\)，\(A_{2}(2,0,0)\)，\(B_{1}(0,1,0)\)，\(B_{2}(0,2,0)\)，\(C_{1}(0,0,1)\)，\(C_{2}(0,0,2)\) 这6个点中随机选取3个点，将这3个点及原点 \(O\) 两两相连构成一个“立体”，记该“立体”的体积为随机变量 \(V\) （如果选取的3个点与原点在同一个平面内，此时“立体”的体积 \(V = 0\)）.
\begin{enumerate}
\item 求 \(V = 0\) 的概率；
\item 求 \(V\) 的分布列及数学期望 \(E(V)\).
\end{enumerate}

\bitmapfigure[max width=0.34\linewidth]{img/2012/jiangxi_science/q19_fig1.png}
\end{problem}

\begin{solution}
\begin{enumerate}
\item 从 \(6\) 个点中任选 \(3\) 个点，共有
\[
\mathrm C_6^3=20
\]
种等可能的选法. 若 \(V\neq0\)，三个所选点必须分别位于 \(x\)、\(y\)、\(z\) 轴上；每条轴上有 \(2\) 个点，故这样的选法有
\[
2\times2\times2=8
\]
种. 因此
\[
P(V=0)=1-\frac8{20}=\frac35
\]

\item 当三个点分别在三条坐标轴上时，设它们到原点的距离分别为 \(u\)，\(v\)，\(w\)，其中 \(u\)，\(v\)，\(w\in\{1,2\}\). 由棱锥体积公式，
\[
V=\frac16uvw
\]
在 \(8\) 种等可能的轴上取点方式中，\((u,v,w)\) 的乘积依次按
\(1\)，\(\ 2\)，\(\ 2\)，\(\ 2\)，\(\ 4\)，\(\ 4\)，\(\ 4\)，\(\ 8\)
出现，所以分布列为

\begin{center}
\begin{tabular}{c|ccccc}
\(V\) & \(0\) & \(\frac{1}{6}\) & \(\frac{1}{3}\) & \(\frac{2}{3}\) & \(\frac{4}{3}\) \\
\hline
\(P\) & \(\frac{3}{5}\) & \(\frac{1}{20}\) & \(\frac{3}{20}\) & \(\frac{3}{20}\) & \(\frac{1}{20}\)
\end{tabular}
\end{center}

并且
\[
E(V)=0\cdot\frac35+\frac16\cdot\frac1{20}
+\frac13\cdot\frac3{20}
+\frac23\cdot\frac3{20}
+\frac43\cdot\frac1{20}
=\frac1{120}+\frac1{20}+\frac1{10}+\frac1{15}
=\frac9{40}
\]
\end{enumerate}
\end{solution}

\begin{problem}
在三棱柱 \(ABC - A_{1}B_{1}C_{1}\) 中，已知 \(AB = AC = AA_{1} = \sqrt{5}\)，\(BC = 4\)，点 \(A_{1}\) 在底面 \(ABC\) 的射影是线段 \(BC\) 的中点 \(O\).
\begin{enumerate}
\item 证明：在侧棱 \(AA_{1}\) 上存在一点 \(E\)，使得 \(OE \perp \) 平面 \(BB_{1}C_{1}C\)，并求出 \(AE\) 的长；
\item 求平面 \(A_{1}B_{1}C\) 与平面 \(BB_{1}C_{1}C\) 夹角的余弦值.
\end{enumerate}

\bitmapfigure[max width=0.42\linewidth]{img/2012/jiangxi_science/q20_fig1.png}
\end{problem}

\begin{solution}
\begin{enumerate}
\item 因为 \(O\) 是 \(BC\) 的中点，所以 \(BO=CO=2\). 又 \(AB=AC\)，故 \(AO\perp BC\)，从而
\[
AO=\sqrt{AB^2-BO^2}=\sqrt{5-4}=1
\]
建立空间直角坐标系，取
\(O=(0,0,0)\)，\(B=(-2,0,0)\)，\(C=(2,0,0)\)，\(A=(0,1,0)\).
由 \(A_1\) 在底面上的射影为 \(O\)，设 \(A_1=(0,0,h)\). 由 \(AA_1=\sqrt5\) 得
\(1+h^2=5\)，\(h=2\).
棱柱的侧棱向量为
\[
\overrightarrow{AA_1}=(0,-1,2)
\]
所以
\(B_1=(-2,-1,2)\)，\(C_1=(2,-1,2)\).
平面 \(BB_1C_1C\) 内有两个不共线方向向量
\(\overrightarrow{BC}=(4,0,0)\)，\(\overrightarrow{BB_1}=(0,-1,2)\).
设 \(E\) 将 \(AA_1\) 按比例 \(s\) 内分，即
\(E=A+s\overrightarrow{AA_1}=(0,1-s,2s)\)，\(0\leqslant s\leqslant1\).
因为 \(\overrightarrow{OE}\cdot\overrightarrow{BC}=0\) 恒成立，要使 \(OE\) 垂直于该平面，还需
\[
\overrightarrow{OE}\cdot\overrightarrow{BB_1}
=(0,1-s,2s)\cdot(0,-1,2)
=-(1-s)+4s=0
\]
解得 \(s=\dfrac15\)，确实属于 \([0,1]\)，所以这样的点存在，且
\[
AE=s\cdot AA_1=\frac15\sqrt5=\frac{\sqrt5}{5}
\]

\item 取平面 \(A_1B_1C\) 内的两个方向向量
\(\overrightarrow{A_1B_1}=(-2,-1,0)\)，\(\overrightarrow{A_1C}=(2,0,-2)\).
其法向量可取
\[
\bs{n}_1=\overrightarrow{A_1B_1}\times\overrightarrow{A_1C}
=(2,-4,2)\sim(1,-2,1)
\]
平面 \(BB_1C_1C\) 的法向量由
\(\overrightarrow{BC}=(4,0,0)\)，\(\overrightarrow{BB_1}=(0,-1,2)\)
给出，可取
\[
\bs{n}_2=\overrightarrow{BC}\times\overrightarrow{BB_1}
=(0,-8,-4)\sim(0,2,1)
\]
两平面夹角取锐角，故其余弦为
\[
\cos\varphi
=\frac{|\bs{n}_1\cdot\bs{n}_2|}
{|\bs{n}_1|\,|\bs{n}_2|}
=\frac{|1\cdot0+(-2)\cdot2+1\cdot1|}
{\sqrt6\sqrt5}
=\frac3{\sqrt{30}}
=\frac{\sqrt{30}}{10}
\]
\end{enumerate}
\end{solution}

\begin{problem}
已知三点 \(O(0,0)\)，\(A(-2,1)\)，\(B(2,1)\)，曲线 \(C\) 上任意一点 \(M(x,y)\) 满足 \(\left|\overrightarrow{MA} + \overrightarrow{MB}\right| = \overrightarrow{OM} \cdot \left(\overrightarrow{OA} + \overrightarrow{OB}\right) + 2\).
\begin{enumerate}
\item 求曲线 \(C\) 的方程；
\item 动点 \(Q(x_0, y_0)\) （\(-2 < x_0 < 2\)）在曲线 \(C\) 上，曲线 \(C\) 在点 \(Q\) 处的切线为 \(l\). 问：是否存在定点 \(P(0, t)\) （\(t < 0\)），使得 \(l\) 与 \(PA\)，\(PB\) 都相交，交点分别为 \(D\)，\(E\)，且 \(\triangle QAB\) 与 \(\triangle PDE\) 的面积之比是常数？若存在，求 \(t\) 的值；若不存在，请说明理由.
\end{enumerate}
\end{problem}

\begin{solution}
\begin{enumerate}
\item 由
\(\overrightarrow{MA}=(-2-x,1-y)\)，\(\overrightarrow{MB}=(2-x,1-y)\)
得
\[
\left|\overrightarrow{MA}+\overrightarrow{MB}\right|
=\sqrt{(-2x)^2+(2-2y)^2}
=2\sqrt{x^2+(1-y)^2}
\]
又
\(\overrightarrow{OA}+\overrightarrow{OB}=(0,2)\)，\(\overrightarrow{OM}\cdot(\overrightarrow{OA}+\overrightarrow{OB})=2y\).
所以题设方程为
\[
2\sqrt{x^2+(1-y)^2}=2y+2
\]
等式左边非负，故等式成立时 \(y+1\geqslant0\)，可以平方，得
\[
x^2+(1-y)^2=(y+1)^2
\]
即
\[
x^2=4y
\]
反之，若 \(x^2=4y\)，则 \(y\geqslant0\)，代回上式两边均为
\(2(y+1)\)，故曲线方程为 \(x^2=4y\).

\item 设 \(Q=(q,q^2/4)\)，其中 \(-2<q<2\)，并记
\(u=\frac{q^2}{4}\in[0,1)\)，\(P=(0,t)=(0,-w)\)，\(w>0\).
抛物线在 \(Q\) 处的切线为
\(l:\)，\(y=\frac q2x-\frac{q^2}{4}=\frac q2x-u\).
用参数表示 \(D\)，\(E\)：
\(D=P+s(A-P)\)，\(E=P+r(B-P)\).
由于 \(A=(-2,1)\)、\(B=(2,1)\)，有
\(D=(-2s,-w+s(1+w))\)，\(E=(2r,-w+r(1+w))\).
将 \(D\)、\(E\) 分别代入切线方程，得到
\(s=\frac{w-u}{1+w+q}\)，\(r=\frac{w-u}{1+w-q}\).
取 \(q\) 足够接近 \(0\)，可使 \(w-u>0\) 且两个分母均为正；若面积之比对所有动点 \(Q\) 为常数，则在这段 \(q\) 上也必须为常数，此时 \(s\)，\(r>0\).
三角形 \(QAB\) 和 \(PAB\) 的底边都是 \(AB=4\)，故
\([QAB]=2(1-u)\)，\([PAB]=2(1+w)\).
又 \(D\)、\(E\) 分别是从 \(P\) 指向 \(A\)、\(B\) 的有向比例为 \(s\)，\(r\) 的点，所以
\[
[PDE]=sr[PAB]=2sr(1+w)
\]
因此面积之比为
\[
R(u)=\frac{[QAB]}{[PDE]}
=\frac{(1-u)\bigl((1+w)^2-4u\bigr)}
{(1+w)(w-u)^2}
\]
若该比值与动点 \(Q\) 无关，则上式关于 \(u\) 的分子、分母必须成比例. 比较 \(u^2\)、\(u\) 和常数项的系数，分别得到
\((1+w)R=4\)，\((1+w)^2=4w^2\)，\((1+w)^2+4=8w\).
因 \(w>0\)，第二个等式给出 \(1+w=2w\)，故 \(w=1\)，从而
\[
t=-w=-1
\]
此时
\(s=\frac{1-u}{2+q}=\frac{2-q}{4}\)，\(r=\frac{1-u}{2-q}=\frac{2+q}{4}\)
且 \(-2<q<2\) 保证 \(0<s\)，\(r<1\)，所以 \(l\) 确实分别与线段 \(PA\)、\(PB\) 相交. 代入面积比得
\[
R(u)=\frac{(1-u)\,4(1-u)}{2(1-u)^2}=2
\]
因此存在定点 \(P(0,-1)\)，且面积之比为 \(2\).
\end{enumerate}
\end{solution}

\begin{problem}
若函数 \(h(x)\) 满足
\begin{circlelist}
\item \(h(0)=1\)，\(h(1)=0\)；
\item 对任意 \(a \in [0,1]\)，有 \(h(h(a)) = a\)；
\item 在 \((0,1)\) 上单调递减.
\end{circlelist}
则称 \(h(x)\) 为补函数. 已知函数 \(h(x) = \left(\frac{1 - x^{p}}{1 + \lambda x^{p}}\right)^{\frac{1}{p}}\)（\(\lambda > -1\)，\(p > 0\)）.
\begin{enumerate}
\item 判断函数 \(h(x)\) 是否为补函数，并证明你的结论；
\item 若存在 \(m \in [0,1]\)，使得 \(h(m) = m\)，称 \(m\) 是函数 \(h(x)\) 的中介元. 记 \(p = \frac{1}{n} \) （\(n \in \mathbf{N}_+\)）时 \(h(x)\) 的中介元为 \(x_n\). 且 \(S_n = \sum_{i=1}^{n} x_i\)，若对任意的 \(n \in \mathbf{N}_+\)，都有 \(S_n < \frac{1}{2}\)，求 \(\lambda\) 的取值范围；
\item 当 \(\lambda = 0\)，\(x \in (0,1)\) 时，函数 \(y = h(x)\) 的图象总在直线 \(y = 1 - x\) 的上方，求 \(p\) 的取值范围.
\end{enumerate}
\end{problem}

\begin{solution}
\begin{enumerate}
\item 对 \(x\in[0,1]\)，有 \(1+\lambda x^p>0\)，且
\[
0\leqslant\frac{1-x^p}{1+\lambda x^p}\leqslant1
\]
其中右侧不等式等价于
\(-(1+\lambda)x^p\leqslant0\). 因此 \(h(x)\) 在 \([0,1]\) 上有定义，且值域包含于 \([0,1]\). 直接代入得
\(h(0)=1\)，\(h(1)=0\).
令
\[
g(x)=h(x)^p=\frac{1-x^p}{1+\lambda x^p}
\]
在 \(0<x<1\) 上，
\[
g'(x)
=-\frac{p(1+\lambda)x^{p-1}}{(1+\lambda x^p)^2}<0
\]
由于 \(p\) 次方根函数递增，\(h(x)\) 在 \((0,1)\) 上单调递减.

再令 \(y=h(x)\)，则
\[
y^p=\frac{1-x^p}{1+\lambda x^p}
\]
于是
\(1-y^p=\frac{(1+\lambda)x^p}{1+\lambda x^p}\)，\(1+\lambda y^p=\frac{1+\lambda}{1+\lambda x^p}\)
从而
\[
\frac{1-y^p}{1+\lambda y^p}=x^p
\]
两边取正的 \(p\) 次方根，得 \(h(y)=x\)，即
\[
h(h(x))=x
\]
所以 \(h(x)\) 是补函数.

\item 设 \(h(m)=m\)，令 \(u=m^p\). 由定义有
\[
u=\frac{1-u}{1+\lambda u}
\]
即
\[
\lambda u^2+2u-1=0
\]
在 \(0\leqslant u\leqslant1\) 内的唯一解为
\[
u=\frac1{1+\sqrt{1+\lambda}}
\]
记
\[
r=\frac1{1+\sqrt{1+\lambda}}
\]
则当 \(p=\dfrac1n\) 时，中介元满足
\[
x_n=r^n
\]
从而
\[
S_n=\sum_{i=1}^nr^i
\]
由于 \(0<r<1\)，\(S_n\) 随 \(n\) 增大，且
\[
\lim_{n\to\infty}S_n=\frac r{1-r}
\]
所有正整数 \(n\) 均满足 \(S_n<\dfrac12\)，当且仅当
\[
\frac r{1-r}\leqslant\frac12
\]
其中等号时每个有限项仍严格小于极限. 于是 \(r\leqslant\dfrac13\)，即
\[
\frac1{1+\sqrt{1+\lambda}}\leqslant\frac13
\Longleftrightarrow \sqrt{1+\lambda}\geqslant2
\Longleftrightarrow \lambda\geqslant3
\]
故 \(\lambda\in[3,+\infty)\).

\item 当 \(\lambda=0\) 时，
\[
h(x)=(1-x^p)^{1/p}
\]
因为 \(p\) 次方根函数递增，\(h(x)>1-x\) 等价于
\[
1-x^p>(1-x)^p
\]
即
\[
x^p+(1-x)^p<1
\]
当 \(p>1\) 且 \(0<x<1\) 时，\(x^p<x\)、\((1-x)^p<1-x\)，所以该不等式成立；当 \(0<p<1\) 时两项不等式方向相反，左侧大于 \(1\)；当 \(p=1\) 时左侧等于 \(1\). 因此所求范围为
\[
p\in(1,+\infty)
\]
\end{enumerate}
\end{solution}
